\section{237塔巡检：完整联合决策模型}
\noindent 本模型在同一个可行域内联合选择任务、航次、车辆道路、停靠时段、无人机连续航迹与姿态、拍摄参数、电池分配和充电。日前规划求一次完整班次；执行中从真实状态重新求同一问题。下列定义针对参考场景，不将模板几何或设备假设称为实测标定。
\begin{center}\small
\begin{tabular}{p{.26\linewidth}p{.36\linewidth}p{.28\linewidth}}\toprule
项目 & 当前完整范围 & 数据性质\\\midrule
区域 & Rommerskirchen；237塔 & OSM杆塔经纬度\\
观测要求 & 2133项；每塔9项 & 文献角色＋几何模板\\
道路 & 13936节点、20455有向弧 & 保留全部源弧与源权限\\
停车候选 & 104个；允许重访 & 模型候选，非许可证明\\
地形 & 1953采样点、3886三角形 & 公开DEM插值\\
资源 & 2车、4机、8对电池 & 每车2机、4对电池\\
共享能力 & 每车1起降位、1充电通道 & 电源1000W/车\\
班次 & 28800s；8h & 任务结束时刻不随重规划后移\\
优先目标 & 完整塔数、返库时间、机载Wh & 依次优化，不混合单位\\\bottomrule
\end{tabular}\end{center}
\noindent 原始场景保留不改。经纬度为EPSG:4326，平面为EPSG:32632；计算前减去固定原点。高程统一使用源DEM基准，原始垂直基准未独立核验。变量采用m、s、W、Wh及rad；同一个时刻用于车辆、飞机、天气和电池。
\begin{center}\small
\begin{tabular}{p{.29\linewidth}p{.61\linewidth}}\toprule
层次 & 本稿采用的精度\\\midrule
物理模型 & 连续时间、非线性、混合离散决策；对全时段施加约束\\
无人机动力学 & 位置三阶导数、力平衡、机体姿态与角速度；非电机级六自由度控制器\\
车辆模型 & 全道路的连续Ackermann运动；地形抬升，非悬架/轮胎模型\\
电池模型 & 物理电池对守恒、分段充电、复用和预留；不优化老化\\
数值实现 & 尚未形成237塔完整SCIP转录；历史9任务结果单独保留\\\bottomrule
\end{tabular}\end{center}
\clearpage

